% Section module for the research paper.  Audit of the pinned arithmetic source,
% not an arithmetic certificate and not a compilation report.
\section{Audit of the arithmetic source and the remaining dependencies}
\label{sec:arithmetic-source-audit}

\subsection{Scope and conclusion}

The source examined here is \texttt{quartic\_12504.tex} at commit
\texttt{28e87b1b62af0effd954605af7f70b2699dd5672}, with SHA--256
\begin{center}\footnotesize\ttfamily
38d458ba4178083766c5ae56a894da595fe27170c1ef07f6b49750fc5edf7c37.
\end{center}
The line numbers below refer to this fixed file.  Labels are also given
because subsequent edits can change line numbers.  We examined the
arithmetic interpretation of the released reciprocal comparison source,
its principal local arguments, its final height comparisons, and selected
substantial inputs to completeness.  We did not independently verify the
entire new restricted converse argument or every application of the
arithmetic-geometric height theorems.

\emph{No rational counterexample to the source system, and no invalid
step in the elementary local and numerical deductions listed below, was
found in this audit.}  This statement is substantially weaker than an
independent proof of the asserted Diophantine definition.  The exact
polynomial transformations elsewhere in this paper preserve the projected
rational zero set.  Identifying that set with $\mathbb Z$ still requires
the arithmetic dependencies isolated here.

Two dependencies deserve particular emphasis.  The pole predicate's
one-sided obstruction has an elementary valuation proof, but its
completeness uses a new restricted $2$-converse.  The asserted
five-contact height theorem has a lengthy geometric and modularity proof;
a calculation of its displayed constant is not a proof of that theorem.
Neither dependency is discharged by checking the polynomial compiler.

\subsection{The local obstruction and norm argument}

For $b\in\mathbb Q$, put $B=17b$ and let $\Psi(b)$ assert the existence
of $t,U,e,p,q,s,h\in\mathbb Q$ satisfying
\[
 U(1+bt^2)=b(1-bt^2),\qquad p^2=B(1-e)(1+3e),\qquad
 h=e(B-U^2),\qquad Bq^2=(s^2-Bh)(s^2+3Bh).
\]
The relevant source statements are
\texttt{hyp:o8-pole} and \texttt{lem:original-seven-parity-soundness}
(lines 7532--7597).

\paragraph{Checked one-sided statement.}
If $b\ne0$ and $\Psi(b)$, then $v_\ell(b)$ is not negative and odd
for any prime $\ell\notin\{2,3,17\}$.  Indeed, suppose
$k=v_\ell(b)<0$ is odd.  The first equation rules out $1+bt^2=0$.
The valuations of $1-bt^2$ and $1+bt^2$ agree, so $v_\ell(U)=k$;
this also holds when $t=0$.  Consequently
$v_\ell(B-U^2)=2k$.  The second equation rules out $e=0$ and forces
$v_\ell(e)=0$: a nonzero valuation would make its right side have
odd valuation.  Thus $h\ne0$ and $v_\ell(Bh)=3k$ is odd.
For any $s$, the right side of the final equation is nonzero and has
even valuation; the two terms in each factor have unequal valuation
when $s\ne0$, and the assertion is immediate when $s=0$.
The left side is either zero or has odd valuation, a contradiction.
This proof includes $e=1$, $e=-1/3$, and $p=0$.  No omitted inverse
condition on a witness is needed.

\paragraph{Checked norm propagation.}
The binary operation used by the compiler is
$N_r(a,b)=a^2+rb^2$.  If $R=v_\ell(r)$ is odd, terms with both inputs
nonzero have distinct valuation parity.  Hence
\[
 v_\ell N_r(a,b)=\min\{2v_\ell(a),R+2v_\ell(b)\}.
\]
In particular a tree containing the constant leaf $1$ cannot have zero
root at such a prime.  If the constant is on the all-left path and
$R<0$, the root $W$ has $v_\ell(W)\le0$; then $rW^2$ has negative
odd valuation, contradicting the pole obstruction.  Thus $R>0$.
The same obstruction then implies
$R+2v_\ell(W)>0$.  Descending the preceding minimum formula shows
that a leaf $X/r^e$ has valuation greater than $-R$.  Therefore
\[
 v_\ell(X)>(e-1)R\ge e-1,
 \qquad\text{and hence}\qquad v_\ell(X)\ge e.
\]
This verifies the strict inequality and integer rounding used in the
mixed-precision batches; it also explains why a global condition
$W\ne0$ is unnecessary for soundness.  See lines 7789--7804 and
8080--8120 of the source.

The finite squareclass correction in lines 7755--7804 is separate from
this soundness argument.  Its extra factors have support in
$\{2,3,5,7,11,13,17\}$, already excluded from the relevant contacts.
Completeness chooses the root to be a unit at primes dividing the raw
integer or the correction.  This condition is needed: multiplying an
arbitrary rational square by an integer does not by itself guarantee a
square reduced denominator.

\subsection{The reciprocal logarithm, common point, and final estimates}

\paragraph{Checked formal approximation.}
For $E:y^2=x^3-Dx$, $D=11-6\sqrt2$, set $f=1/x$, $z=-x/y$,
$a_0=1$, $a_1=-2/5$, and
$a_j=(4j-5)a_{j-1}/(4j+1)$.  The relations
\[
 f=z^2(1-Df^2),\qquad
 4w(1-w)G'(w)+(1+w)G(w)=1-w,
 \quad G(w)=\sum_{j\ge0}a_jw^j,
\]
give $\log_E(z)=zG(Df^2)$.  At a good split prime $p>76176$,
with $r=v_p(z)\ge1$, one has $v_p(f)=2r$.
For $e\in\{16,64\}$, $m=e/4$, and
$L_e=z\sum_{j=0}^{m-1}a_j(Df^2)^j$, the source bounds the entire
tail, including indices whose denominators are divisible by $p$:
\[
 v_p(\log_E(z)-L_e)\ge(4m+1)r,
 \qquad L_e/\log_E(z)\in1+p^e\mathbb Z_p.
\]
The factorial denominator bound used for the tail is sufficient.
For $4j+1\ge p$ one has $j\ge19044$ and
$v_p(a_j)\ge-(4j+1)/(p-1)\ge-j$, which gives the required bound
and convergence.  Truncating before the first denominator divisible by
$p$ alone would not justify the infinite-tail estimate; the source does
not make that shortcut.  See \texttt{lem:quartic-reciprocal-log},
lines 8753--8816.

The separate global nonvanishing argument is also valid for the
truncations used here.  In $\mathbb Q(\sqrt2)$ the valuation extending
$v_2$ has value group $\frac12\mathbb Z$, $v(D)=0$, and $v(a_j)=1$
for $j\ge1$.  If $v(f)\ge0$, the constant term of $G_J(Df^2)$
is uniquely least; if $v(f)<0$ and $J\ge2$, its last term is uniquely
least.  Thus the logarithm inverses introduce no exceptional rational
point.  See \texttt{lem:quartic-reciprocal-nonzero}, lines 8818--8842.

\paragraph{Checked compatibility of the common auxiliary point.}
In the completeness construction the private congruences are
\[
 h_p\equiv0\pmod{d_{p,1}},\quad
 h_p\equiv4\pmod{d_{p,2}},\quad
 h_p\equiv4\pmod{p^{64+v_p(j)}},\qquad
 j=\operatorname{lcm}_{p}d_{p,2}.
\]
The conditions $\gcd(d_{p,1},d_{p,2})=4$ and
$p\nmid d_{p,1}d_{p,2}$ suffice for these congruences at each
individual $p$.  The construction does not need pairwise coprimality
between group orders belonging to different primes.  A selected prime
may divide $j$ because of another group order; the term $v_p(j)$
already accounts for this possibility.  Avoiding the finitely many
forbidden values of $h_p$ preserves the point-chart guards.
The same verification does not authorize replacing all $h_p$ by a single
shared $h$: that would require compatibility between congruences with
moduli $d_{p,1}$ and $d_{q,2}$ for different primes $p,q$, and the source
provides no such cross-prime hypothesis.  Sharing this witness is therefore
a blocked optimization unless an additional prime-pattern theorem or a
different construction is supplied.  This checks the stated use of the
ordinary-prime pattern, rather than independently reproving the full
prime-pattern theorem; see
\texttt{subsec:o8-common-J-proof}, lines 8394--8468, and lines
8964--8984.

\paragraph{Checked final numerical implication.}
Assume the stated height and radical-divisibility inputs, and put
$a=|s|$, $B=\max(|n|,|s|,|t|,|z|)$, $H=H_*$.
The anchor gives $a>(20H)^2$.  If the integer Lucas residual
\[
 \Delta=t^2-s(s-2)tz+(s-2)^2z^2-1
\]
is nonzero, the source obtains
$B^2-6<7H|\Delta|^{3/8}<14HB^{3/2}$.
Dividing by $B^{3/2}$, with $B\ge1$, yields
$\sqrt B<14H+6\le20H$, a contradiction.  Thus $\Delta=0$.
The elementary Lucas-unit argument then gives $\log B>a-3$.

If $\beta\ne n$, let $N$ be the absolute numerator of
$\beta-n$ in lowest terms.  The precision-$16$ argument gives
\[
 B^2<7HN^{3/2}+6,\qquad N\le2e^{h(\beta)}B.
\]
Consequently
$\log B<2\log(32H)+3h(\beta)$.
The scalar estimate $h(\beta)\le H\sqrt a<a/20$ and
$\log(32H)<a/20$ make the right side less than $a/4<a-3$.
There is no loss of the exponent when the beta precision changes
from $64$ to $16$.  Finally $\beta=n$ and
$X^2=\beta-2\in\mathbb Z$ imply $X\in\mathbb Z$ by coprime
numerator and denominator.  See \texttt{thm:quartic-source-arithmetic},
lines 8867--8985.  This is a conditional verification of the final
implication, not a verification of its height hypotheses.

\subsection{The two substantial arithmetic dependencies}

\paragraph{Restricted converse and pole completeness.}
The exact converse needed is
\[
 \operatorname{corank}_{\mathbb Z_2}
  \operatorname{Sel}_{2^\infty}(E_\ell/\mathbb Q)=1
 \quad\Longrightarrow\quad
 \operatorname{ord}_{s=1}L(E_\ell,s)=1,
 \qquad E_\ell:y^2=x(x-\ell)(x+3\ell),
\]
for positive squarefree $\ell\equiv11\pmod{12}$.
It is \texttt{hyp:RC} at lines 1016--1028.  The manuscript supplies
its own long argument in Sections \texttt{sec:rc},
\texttt{sec:arithmetic}, and \texttt{sec:limits} (lines 1771--5755),
with appendices at lines 9589--13699.  It is not an immediate
consequence of ordinary two-descent or of the published forward
analytic-rank-zero-and-one theorems.

The distinction is essential.  A two-Selmer group of dimension three,
in the presence of full rational two-torsion, implies
$r+\dim_{\mathbb F_2}\Sha[2]=1$.  The Cassels--Tate pairing permits
one to infer full $2^\infty$-Selmer corank one, as in
\texttt{lem:corank}; it does not by itself remove a divisible
$\mathbb Q_2/\mathbb Z_2$ summand of $\Sha$ or prove analytic rank one.
The new restricted converse is used precisely at this remaining step.

We checked the algebraic shape of the simultaneous realization and
its reduction to this converse, but have not certified all its local
descent tables or the new converse.  The most important further audit
within the latter is the passage from finite Heegner classes and local
conditions to the uniform valuation bound and the fixed-field
coefficient transfer:
\texttt{ex:limits-uniform-lower-bound} (lines 4920--5217),
\texttt{ex:limits-fixed-field-clearing} (lines 5498--5657), and
\texttt{ex:limits-coefficient-transfer} (lines 5658--5754).
These statements assert uniformity over arbitrarily large active
supports and are part of the new arithmetic argument.  The present
report neither proves nor disproves them.

\paragraph{Five-contact height theorem and its exact constant.}
The required statement is \texttt{hyp:o8-height}
(lines 6666--7520):
\[
 h(a)\le H_*M_5(a)^{24},\qquad \log_2H_*<4783227,
\]
where the contacts are $\infty,0,1,2401,-343/32$ and the product
$M_5(a)$ omits primes at most $76176$.
The displayed bound on $H_*$ is needed for the released anchor;
an unspecified larger constant cannot simply be substituted.
The source derives it using the discriminant-$210$ Shimura curve,
a degree-at-most-$576$ normal fiber field, a comparison with stable
Faltings height, and separate CM and non-CM estimates.

Two plausible objections were checked against primary sources and
should not be presented as gaps.  Yuan's Theorem 1.1 and Section 2.1
of \emph{Explicit Kodaira--Spencer map over Shimura Curves}
(arXiv:2205.11334v2, pp. 2--6) do provide the stated maximal-stack
regularity and semistability and the isometry whose image is
$d_B\omega^{\otimes2}$.  Pan's Theorem 1.0.4 of
\emph{The Fontaine--Mazur conjecture in the residually reducible case}
(arXiv:1901.07166v2, p. 5) has no extra residual-image restriction
at $p=5$, provided the representation is irreducible, odd, finitely
ramified, potentially semistable, and regular.  Thus residual
reducibility alone does not invalidate the source's use of Pan.

What remains to be checked is the complete chain of hypotheses in the
source's application: the canonical normalizer action and constant
fiber torsor, the integral height comparison with its normalization,
the scalar cocycle splitting and ramification bounds, and the
modular-Jacobian transfer preserving the displayed constant.
The principal non-CM passage is
\texttt{sec:h5-level-repair}, lines 6904--7345; the CM estimate follows
at lines 7349--7500.  These are substantive mathematical arguments,
not consequences of the finite numerical certificate.  Source comments
marking a repair as awaiting review record a review status; they are
not, by themselves, a proof of an error.

\subsection{Consequences for the optimization claims}

The unconditional conclusions established by an exact transformation
certificate have the form
\[
 \{X:\exists\mathbf y\ F(X,\mathbf y)=0\}
 =\{X:\exists\mathbf z\ F_{\rm new}(X,\mathbf z)=0\}.
\]
They remain meaningful independently of the source's main arithmetic
claim.  Until the dependencies above have been independently verified,
the further conclusion that either set equals $\mathbb Z$ must be
stated as conditional on the arithmetic theorem of the source.
A counterexample to an individual intermediate lemma would require
tracing its use and possible repair before concluding that the released
polynomial fails; no such counterexample is supplied by this audit.
